% Lista MANUAL em ordem alfabética, sem numeração, conforme MLA 9/RDQ.
% Inclua apenas as obras realmente citadas e confira os dados da edição usada.
% Títulos significativos em maiúsculas iniciais; livros em itálico;
% artigos/capítulos entre aspas; não informar cidade para livros.
% Referências abreviadas ficam NAS NOTAS, não em chamadas autor-data.

\referencia{Agência Câmara de Notícias. \enquote{Joice Hasselmann Diz que Deixa a Liderança no Congresso com Sentimento de Dever Cumprido}. \textit{Câmara dos Deputados}, 17 out. 2019. Disponível em: \url{https://www.camara.leg.br/noticias/601207-joice-hasselmann-diz-que-deixa-a-lideranca-no-congresso-com-sentimento-de-dever-cumprido/}. Acesso em: 21 nov. 2026.}

\referencia{Barros, Douglas Rodrigues. \enquote{Contra Nego Bispo}. \textit{Blog da Boitempo}, 21 jan. 2026. Disponível em: \url{https://blogdaboitempo.com.br/2026/01/21/contra-nego-bispo/}. Acesso em: 20 nov. 2026.}

\referencia{Carneiro da Cunha, Manuela. \enquote{\enquote{Cultura} e cultura: conhecimentos tradicionais e direitos intelectuais}. \textit{Cultura com Aspas e Outros Ensaios}. Ubu Editora, 2017.}

\referencia{Foucault, Michel. \enquote{As heterotopias}. \textit{O corpo utópico, as heterotopias}, n-1 edições, 2013.}

\referencia{Foucault, Michel. \textit{As Palavras e as coisas}, Martins Fontes, 1999.}

\referencia{Haraway, Donna. \enquote{Pensamento tentacular: Antropoceno, Capitaloceno, Chthuluceno}. \textit{Ficar com o problema: Fazer parentes no Chthluceno}, n-1 edições, 2023. p 59-116.}

\referencia{Krenak, Ailton. \textit{Futuro Ancestral}. Companhia das Letras, 2022.}

\referencia{Latour, Bruno. \enquote{How to Talk About the Body? The Normative Dimension of Science Studies}. \textit{Body \& Society}, vol.~10, n.~2--3, 2004, pp.~205--229. \url{https://doi.org/10.1177/1357034X04042943}.}

\referencia{Latour, Bruno. \textit{Jamais Fomos Modernos: Ensaio de Antropologia Simétrica}. Traduzido por Carlos Irineu da Costa, Editora 34, 2019.}

\referencia{Mauss, Marcel; Beuchat, Henri. \enquote{Ensaio sobre as variações sazonais das sociedades esquimós}. \textit{Sociologia e Antropologia}, escrito por Marcel Mauss, Cosac Naify, 2015. p. 423-501.}

\referencia{R7. \enquote{Ex-Líder do Governo Bolsonaro, Joice Hasselmann Declara Voto em Lula no 2º Turno}. \textit{R7 Notícias}, 23 out. 2026. Disponível em: \url{https://noticias.r7.com/eleicoes/2026/ex-lider-do-governo-bolsonaro-joice-hasselmann-declara-voto-em-lula-no-2-turno-23102026/}. Acesso em: 21 nov. 2026.}

\referencia{Santos, Antônio Bispo dos. \enquote{Dois e Dois São Dois: Renato Noguera e Nêgo Bispo}. Entrevista concedida a Renato Noguera. \textit{Amarello}, n.~40, 2022. Disponível em: \url{https://amarello.com.br/2022/04/cultura/dois-e-dois-sao-dois-renato-noguera-e-nego-bispo/}. Acesso em: 20 nov. 2026.}

\referencia{Santos, Antônio Bispo dos. \enquote{Entrevista: Nego Bispo}. Entrevista concedida a Gabriel Galo em 17 e 18 de junho de 2020. \textit{Papo de Galo}, n.~3, 2020. Disponível em: \url{https://papodegalo.com.br/entrevista-nego-bispo/}. Acesso em: 20 nov. 2026.}

\referencia{Santos, Antônio Bispo dos. \textit{A Terra Dá, a Terra Quer}. Piseagrama e Ubu Editora, 2023.}

\referencia{Santos, Gabriel. \enquote{Uma Crítica Marxista ao Pensamento Decolonial e a Nego Bispo}. \textit{Esquerda Online}, 10 fev. 2026. Disponível em: \url{https://esquerdaonline.com.br/2026/02/10/uma-critica-marxista-ao-pensamento-decolonial-e-a-nego-bispo/}. Acesso em: 20 nov. 2026.}